\documentclass[12pt]{article}
\input{bayesuvius.sty}
\begin{document}


Below, $\indi()$ 
will denote the indicator
or truth function.
$\indi(S)$ equal 1 if $S$ is true,
and 0 if $S$ is false.
For example, $\indi(x\geq 0)$
equals 0 if $x<0$ and 1 otherwise.

{\bf Standardized Precipitation Index} (SPI) ranges from $-4$ (least rain) to $4$ (most rain).



$I_D=$ {\bf Impact } (i.e., {\bf severity}) of drought is defined as minus SPI (higher rain,
higher SPI, less 
severe drought)

$ P(I_D)=$ likelihood, 
probability  that a drought of severity $I_D$  will occur in the near future (weeks to months), based on forecasts.

$\mu=$ ensemble member index

$M=$ number of ensemble members, 51 for $ECMWF$ (European Centre for Medium-Range Weather Forecasts)


\beqa
Risk_D &=& E[\rvI_D] = \frac{1}{M} \sum_{\mu=1}^{M} I_{D,\mu}
\\
&=&\sum_{I_D} P(I_D) I_D
\eeqa

A {\bf risk matrix} is
a  $4\times 4$ matrix
with $X$ axis 
equal to impact ($I$)
and $Y$ axis equal to likelihood
$P$, with 
$I$ binned into four values
and $P(I)$ binned into 4 values.
The matrix elements equal
to some arbitrary function 
$f(I, P)$ defined as risk. Usually, $f(I, P)=IP$ but not necessarily. I find 
risk matrices arbitrary
and not very useful. The 
Wikipedia article on risk matrices (Ref.\cite{wiki-risk-matrix}) agrees.


$I_D^*  =$ impact threshold 
   for drought (e.g., $I_D^* =-SPI=1.5$)
   





\beq
P(\rvI_D\geq I_D^*  ) = \frac{1}{M} \sum_{\mu=1}^{M} \indi(I_{D,\mu} \geq I_D^*) 
\eeq

$P^*_D=$ probability threshold for drought

$C_D=$ cost of drought

$L_D=$ loss of drought

\beq
P^*_D = \frac{C_D}{L_D}
\eeq

$T_D=$ trigger for drought


\beq
T_D=
\indi\left[P(\rvI_D\geq I_D^*  ) > P^*_D\right]
\eeq

Let
$N^D$ be the number of times (seasons) 
in which predicted and observed droughts were compared.

 
$N^D_{y|x}=$  number of events in which 
$x\in \bool$ droughts occurred and $y\in \bool$ droughts were predicted

For example, $N^D_{1|0}$ is the number of events 
in which no drought actually occurred but 
it was predicted nonetheless (False Alarm)

for clarity, will omit $D$ superscript from $N^D$ below

\beq
N_{|0} = N_{0|0} + N_{1|0}
\eeq

\beq
N_{|1}= N_{0|1} + N_{1|1}
\eeq

\beq
N = N_{|0}  + N_{|1}
\eeq

\beq
\text{{\bf Hit Rate} (HR)} = P(1|1)
=\frac{N_{1|1}}{N_{|1}}
\eeq

\beq
\text{{\bf False Alarm Rate} (FAR)} = P(1|0)
=\frac{N_{1|0}}{N_{|0}}
\eeq

\beq
\text{{\bf Miss Rate} (MR)} = P(0|1)
=\frac{N_{0|1}}{N_{|1}}
\eeq

\beq
\text{{\bf Specificity}} = P(0|0)
=\frac{N_{0|0}}{N_{|0}}
\eeq




\bibliographystyle{plain}
\bibliography{aa_references}
\end{document}